\section{The Locked \texorpdfstring{\(\Dfour\)}{D4} and \texorpdfstring{\(\Vfour\)}{V4} Subgroup Rows}
\label{sec:d4-v4}

\begin{definition}[Composition and coset convention]
\label{def:composition-convention}
In this section a word \(p=p_1p_2p_3p_4\) is read as the function \(i\mapsto p_i\). Permutations are composed by
\[
  \operatorname{compose}(p,q)=p\circ q,
\]
that is, \(p\) after \(q\). For a subgroup \(H\le \Sfour\), left cosets are written \(gH\) and right cosets are written \(Hg\).
\end{definition}

\begin{definition}[Locked rows]
\label{def:locked-d4-v4-rows}
The locked rows used in this paper are
\[
\Dfour=\{\word{1234},\word{1324},\word{2143},\word{2413},\word{3142},\word{3412},\word{4231},\word{4321}\}
\]
and
\[
\Vfour=\{\word{1234},\word{2143},\word{3412},\word{4321}\}.
\]
\end{definition}

\begin{corollary}[The locked \texorpdfstring{\(\Dfour\)}{D4} and \texorpdfstring{\(\Vfour\)}{V4} rows are not exact poset cones]
\label{cor:d4-v4-not-poset-cones}
The rows \(\Dfour\) and \(\Vfour\) in \cref{def:locked-d4-v4-rows} are not exact poset cones.
\end{corollary}

\begin{proof}
For each unordered pair \(\{a,b\}\) with \(a<b\), the word \(\word{1234}\) places \(a\) before \(b\). The following table gives, for each locked row, a word placing \(b\) before \(a\).
\[
\begin{array}{c|cc}
\{a,b\} & \Dfour\text{ reverse witness} & \Vfour\text{ reverse witness}\\
\hline
\{1,2\} & \word{2143} & \word{2143}\\
\{1,3\} & \word{3142} & \word{3412}\\
\{1,4\} & \word{2413} & \word{3412}\\
\{2,3\} & \word{1324} & \word{3412}\\
\{2,4\} & \word{3142} & \word{3412}\\
\{3,4\} & \word{2143} & \word{2143}
\end{array}
\]
Thus no ordered pair of distinct labels is a common-precedence relation for either locked row. In both cases \(P_X\) is the antichain on \([4]\), and \(\Lin(P_X)=\Sfour\). The locked rows have cardinalities \(8\) and \(4\), respectively, so neither row is equal to its common-precedence closure. By \cref{prop:common-precedence-exactness}, neither row is an exact poset cone.
\end{proof}

\begin{corollary}[The locked rows are subgroup rows with induced coset partitions]
\label{cor:d4-v4-subgroup-cosets}
The locked \(\Dfour\) row is a subgroup row of order 8, so its left and right cosets partition \(\Sfour\) into three blocks. The locked \(\Vfour\) row is a subgroup row of order 4, so its left and right cosets partition \(\Sfour\) into six blocks.
\end{corollary}

\begin{proof}
For the locked \(\Dfour\) row, set
\[
  r=\word{2413},\qquad s=\word{1324}.
\]
Under the convention of \cref{def:composition-convention}, direct composition gives
\[
  r^4=\word{1234},\qquad s^2=\word{1234},\qquad srs=r^{-1}=\word{3142}.
\]
The normal forms generated by \(r\) and \(s\) are
\[
\begin{array}{llll}
r^0=\word{1234}, & r^1=\word{2413}, & r^2=\word{4321}, & r^3=\word{3142},\\
sr^0=\word{1324}, & sr^1=\word{3412}, & sr^2=\word{4231}, & sr^3=\word{2143}.
\end{array}
\]
These are exactly the eight words listed for \(\Dfour\). Hence the locked \(\Dfour\) row is the subgroup \(\langle r,s\rangle\) of order \(8\).

For the locked \(\Vfour\) row, set
\[
  a=\word{2143},\qquad b=\word{3412}.
\]
Direct composition gives
\[
  a^2=\word{1234},\qquad b^2=\word{1234},\qquad ab=ba=\word{4321}.
\]
Therefore
\[
  \Vfour=\{\word{1234},a,b,ab\}
\]
is a subgroup of order \(4\).

Finally, in any finite group, the left cosets of a subgroup partition the group, and the right cosets also partition the group. Since \(|\Sfour|=24\), the number of left or right cosets is the index \(|\Sfour|/|H|\). Thus the locked \(\Dfour\) subgroup has \(24/8=3\) left cosets and \(3\) right cosets, while the locked \(\Vfour\) subgroup has \(24/4=6\) left cosets and \(6\) right cosets.
\end{proof}

\begin{remark}[Subgroup separation already visible]
\label{rem:d4-v4-separation-visible}
\Cref{cor:d4-v4-not-poset-cones,cor:d4-v4-subgroup-cosets} show that subgroup-row status does not imply exact poset-cone status for the selected locked rows. The coset counts are the standard partitions induced by those subgroup rows; they are not used as an independent structural axis.
\end{remark}